\lysh{22268_SOLUTION}{1}
\text{\Large \textbf{ΛΥΣΗ}}

Η εξίσωση (1) είναι της μορφής $\frac{x^2}{a^2} + \frac{y^2}{\beta^2} = 1$, όπου $a^2=9$ και $\beta^2 = 4$. Η μορφή αυτής της εξίσωσης παριστάνει τα σημεία του επιπέδου που βρίσκονται σε έλλειψη με εστίες στον άξονα $x'x$. Οι εστίες έχουν συντεταγμένες $E(\gamma, 0)$, και $E'(-\gamma, 0)$, όπου $\gamma = \sqrt{a^2 - \beta^2} = \sqrt{9-4} = \sqrt{5}$. Η εκκεντρότητα της έλλειψης είναι $\varepsilon = \frac{\gamma}{a} = \frac{\sqrt{5}}{3}$. Το μήκος του μεγάλου άξονα της έλλειψης είναι ίσο με $2a = 2 \cdot 3 = 6$.

α) «Τα σημεία του επιπέδου που επαληθεύουν την εξίσωση (1) βρίσκονται σε μια καμπύλη που ονομάζεται \textbf{έλλειψη}. Οι εστίες της $E$ και $E'$, έχουν συντεταγμένες $E(\sqrt{5}, 0)$ και $E'(-\sqrt{5}, 0)$. Το μήκος του μεγάλου άξονα είναι ίσο με \textbf{6} και η εκκεντρότητα της είναι ίση με $\frac{\sqrt{5}}{3}$».

β)
% TODO: TikZ σχήμα

Η εφαπτόμενη ευθεία σε σημείο με συντεταγμένες $(x_1, y_1)$ της έλλειψης είναι της μορφής
\[
\varepsilon: \frac{x x_1}{9} + \frac{y y_1}{4} = 1 \Leftrightarrow 4 \cdot x x_1 + 9 \cdot y y_1 = 36.
\]
Δίνεται το σημείο επαφής $B(0, -2)$, οπότε αν θέσουμε στην εξίσωση της ευθείας $\varepsilon$ όπου $x_1 = 0$ και $y_1 = -2$ θα έχουμε
\[
\varepsilon: 4 x \cdot 0 + 9 y \cdot (-2) = 36 \Leftrightarrow \varepsilon: -18 y = 36 \Leftrightarrow \varepsilon: y = -2.
\]
\end{lysh}